% Enhanced Literature Review with Full Paper Summaries

In macro and financial economic forecasting, the yield curve has been documented as the leading indicator for economic recessions. This literature review examines the theoretical and empirical foundations of yield curve-based recession forecasting, with particular emphasis on factor decomposition methods, term premium adjustments, and the evolution of predictive modeling techniques.

\citet{estrellaYieldCurvePredictor1996} establish the foundational empirical framework for using the yield curve to predict U.S. recessions. Using data from 1960 to 1995, they demonstrate that the spread between the 10-year Treasury note and the 3-month Treasury bill is a powerful predictor of recessions at horizons of 2 to 6 quarters ahead. Their probit model framework shows that the yield spread significantly outperforms other widely used forecasting measures, including the stock market, the Commerce Department's index of leading economic indicators, and various monetary aggregates. The authors find that a spread greater than 1.2 percentage points corresponds to a recession probability of less than 5 percent, while the probability rises to approximately 25 percent when the spread narrows to zero, and jumps to 50 percent when the curve inverts to a negative 0.8 percentage point spread. The predictive strength stems from the yield curve's role as a reflection of market anticipations regarding future monetary policy easing. When investors expect economic contraction, they anticipate that the Federal Reserve will lower short-term rates, causing the curve to flatten or invert. The paper's out-of-sample forecasting tests demonstrate robust predictive performance across different time periods, establishing the yield spread as one of the most reliable leading indicators available to policymakers and market participants.

\citet{diebold2006forecasting} propose a dynamic Nelson-Siegel framework that revolutionizes yield curve modeling by reducing the entire term structure to three interpretable latent factors: level, slope, and curvature. Building on the classic Nelson-Siegel (1987) exponential components framework, they treat these three factors as time-varying parameters that evolve according to autoregressive processes. The level factor ($\beta_1$) represents long-term interest rates and shifts the entire curve vertically, exhibiting the highest persistence with correlation of 0.99 to empirical measures. The slope factor ($\beta_2$), defined as the negative spread between short-term and long-term rates, shows correlation of 0.99 with the traditional 3-month to 10-year spread and is closely connected to the monetary policy instrument. The curvature factor ($\beta_3$) captures the convexity or "hump" in the yield curve, with peak loading at medium-term maturities around 30 months. While their primary focus is on optimizing the accuracy of yield level forecasts rather than recession prediction, their framework proves particularly valuable for forecasting at horizons of 12 months and beyond, where it substantially outperforms random walk and Fama-Bliss forward rate regression models. The model's 12-month-ahead forecasts achieve significant improvements in root mean squared error at all maturities, often by wide margins. This data-driven factor extraction method provides a parsimonious yet highly effective representation of yield curve dynamics, capturing both cross-sectional fit and time-series evolution without imposing the strict no-arbitrage constraints of affine term structure models. The framework has become widely adopted by central banks globally for fitting and forecasting government bond yields.

\citet{rudebusch2006macroeconomic} provide crucial insights into the dual nature of long-term interest rates by decomposing movements into two distinct components: changes in expectations about future short-term rates and variations in the term premium (the compensation investors demand for holding long-term bonds rather than rolling over short-term securities). Their analysis demonstrates that the macroeconomic implications of these two components can differ significantly, challenging traditional yield curve analysis that treats the spread as a monolithic indicator. Using data from the U.S. term structure and employing various estimation methods including the Kim-Wright model, they show that both monetary policy shocks and technology shocks affect the term premium and output in opposite directions—monetary and technology shocks raise the term premium but decrease output, while demand shocks increase both. The paper documents a significant decline in GDP volatility since 1984 (following McConnell and Perez-Quiros 2000) and an additional sharp drop between 2001 and 2005, suggesting that reduced macroeconomic volatility has contributed to lower term premia. They estimate that reduced volatility and increased Federal Reserve predictability may account for as much as 40 basis points of the decline in long-term relative to short-term yields during the mid-2000s monetary tightening cycle. This is particularly relevant for understanding the atypical 2004-2005 episode when the Federal Reserve increased the federal funds rate from 1 percent to 4.25 percent, yet 10-year Treasury yields fell from 4.7 percent to 4.5 percent—contrary to the typical positive correlation between short-term and long-term rates during tightening cycles. The research highlights that failing to account for term premium dynamics can lead to misinterpretation of yield curve signals, especially during periods of unconventional monetary policy such as quantitative easing or when interest rates approach the zero lower bound.

\citet{bauer2018information} provide a comprehensive empirical assessment of which term spread measures are most reliable for forecasting U.S. recessions and whether adjustments for term premia enhance predictive accuracy. Using monthly data from January 1972 to July 2018 and employing area under the curve (AUC) metrics to measure forecast accuracy, they systematically compare the predictive power of various spread measures for identifying recessions 12 months ahead. Their key finding is that the 10-year minus 3-month Treasury spread emerges as the most reliable predictor, achieving an AUC of approximately 0.88, substantially better than the 0.5 that would correspond to a coin toss. This spread outperforms the more commonly cited 10-year minus 2-year spread ($\text{AUC} \approx 0.82$) and the 10-year minus 1-year spread. Importantly, they find that adjusting the term spread for estimates of the term premium—contrary to some theoretical arguments—does not improve recession forecasting accuracy and may even diminish it. Specifically, when they decompose the term spread into its expectations component and risk premium component using various term premium estimates (including Kim-Wright and Adrian-Crump-Moench models), both components carry similar signals about future recessions, but combining them or using adjusted spreads provides no additional forecast improvement beyond using the unadjusted spread alone. The authors argue that the yield curve requires an actual inversion to reliably signal recession—a merely flat curve is insufficient. Their analysis reveals that every U.S. recession in the preceding 60 years was preceded by a negative term spread, and except for one false positive in the mid-1960s, a negative term spread was always followed by an economic slowdown. They caution, however, that while these empirical correlations are robust, they do not establish causation, leaving open whether yield curve inversions cause recessions (perhaps through credit channel effects as banks face compressed net interest margins) or merely reflect market participants' recession expectations formed on other grounds. This distinction has important implications for policy interpretation during periods of yield curve flattening.

Building upon this established literature, our approach extends the foundational work of \citet{estrellaYieldCurvePredictor1996} in several key directions. First, reflecting the market's collective conviction on future economic growth and monetary policy cycles, we adopt the factor decomposition methodology inspired by \citet{diebold2006forecasting}, but extract level, slope, and curvature factors through Principal Component Analysis rather than parametric Nelson-Siegel estimation, providing a purely data-driven representation unconstrained by functional form assumptions. Within this framework, the level factor (PC1\_Level) represents long-term interest rates and inflation expectations, effectively shifting the entire curve vertically. The slope factor (PC2\_Slope), defined as the spread between long-term and short-term rates, serves as the primary recessionary barometer; a negative slope (inversion) implies that investors anticipate future rate cuts due to economic contraction. Finally, the curvature factor (PC3\_Curvature) captures the "hump" or convexity of the curve, representing the movement of medium-term yields relative to the short and long ends, which is particularly sensitive to market volatility and unconventional policy interventions.

Following the insights of \citet{rudebusch2006macroeconomic}, who highlight that movements in long-term rates represent both changes in expectations about future short rates and variations in the term premium with potentially divergent macroeconomic implications, we incorporate term premium adjustments into our framework. Using the Kim-Wright 10-year term premium estimate (tp\_10y\_kw), we derive an adjusted slope series (slope\_ex\_tp) to discern whether curve inversion is driven by true expectations of monetary easing or a mere compression of the term premium—a distinction that is especially relevant during periods of Quantitative Easing (QE) or the Zero Lower Bound (ZLB). While \citet{bauer2018information} argue that adjusting for term premia does not necessarily improve accuracy in a model-free setting, we embed the term premium and adjusted slope directly into a multivariate framework that includes real-sector indicators, such as unemployment (unrate) and industrial production (indpro\_yoy), to account for sudden exogenous shocks like the COVID-19 pandemic.

This theoretical decomposition is clearly visible in the transition from the pre-crisis to the recovery era shown in Figure \ref{fig:us_yield_curve_examples}. The 2006 average (blue line) exhibits a "high Level, zero Slope" dynamic, where short-term rates ($y_{3m} \approx 5\%$) have converged with long-term rates ($y_{30y} \approx 5\%$). The absence of Slope and Curvature here indicates a restrictive monetary environment where the term premium has eroded, a classic precursor to the 2008 recession. Conversely, the 2010 average (red line) illustrates a fundamental regime shift. While the level at the long end remains elevated, the slope has steepened dramatically ($S_t \gg 0$) as short-term rates approach the zero lower bound. Furthermore, the 2010 curve displays significant curvature (concavity), reflecting the market's nuanced pricing of medium-term recovery risk against a backdrop of aggressive monetary stimulus.

To address the limitations of linear probit models, our approach employs machine learning techniques such as Random Forests and Gradient Boosting, which allow us to capture regime dependency, structural changes, and complex nonlinear interactions that may vary across different economic environments. Finally, we implement a strictly time-respecting, out-of-sample evaluation framework, ensuring that model parameters, classification thresholds, and feature selection are all determined without look-ahead bias, thereby providing a more robust real-time prediction of recessionary risks that would be operationally feasible for policymakers and market participants.

\begin{figure}[H]
    \centering
    \includegraphics[width=0.85\textwidth]{figures/fig_us_yield_curve_0610.pdf}
    \caption{US Yield Curve for selected Maturities for 2006 and 2010 \\Source: FRED}
    \label{fig:us_yield_curve_examples}
\end{figure}